% Source: https://artofproblemsolving.com/community/c3h1573693_area_of_a_triangle
% Stable identity: aops-topic-1573693-q02
The side from $\left(-\dfrac{\pi}{3},\dfrac{\pi}{3}\right)$ to $\left(\dfrac{\pi}{3},\dfrac{\pi}{3}\right)$ is horizontal with length $\dfrac{2\pi}{3}$. The vertex $(0,\pi)$ is at height $\pi-\dfrac{\pi}{3}=\dfrac{2\pi}{3}$ above it (dashed).

\begin{asy}
usepackage("amsmath");
size(5cm);
defaultpen(fontsize(10pt));
import graph;
real p=pi;
pair A=(0,p),B=(-p/3,p/3),C=(p/3,p/3);
draw((-2.6,0)--(2.6,0),gray,Arrow(TeXHead));draw((0,-0.4)--(0,3.8),gray,Arrow(TeXHead));
label("$x$",(2.6,0),E,gray);label("$y$",(0,3.8),N,gray);
fill(A--B--C--cycle,paleyellow);draw(A--B--C--cycle,linewidth(0.9));
dot(A);dot(B);dot(C);
label("$(0,\pi)$",A,NE);label("$\left(-\dfrac{\pi}{3},\dfrac{\pi}{3}\right)$",B,W);label("$\left(\dfrac{\pi}{3},\dfrac{\pi}{3}\right)$",C,E);
draw(A--(0,p/3),dashed);
label("\mbox{}",(min(currentpicture,user=true)+max(currentpicture,user=true))/2,black); // keeps the text after the figure black when LaTeX typesets the labels inline
\end{asy}

So the area is
\[
\dfrac{1}{2}\cdot\dfrac{2\pi}{3}\cdot\dfrac{2\pi}{3}=\boxed{\dfrac{2\pi^2}{9}}.
\]
